\chapter{The two ends of the gauge: what moves a leukocyte}\label{ch:leukocyte}

\section{Why blood first}

The instrument reads whole blood first because it can: leukocytes are the majority of the DNA in a blood draw, their reference methylomes are the best
measured of any human cells, and the known-mixture arrays let the composition be checked against truth. Everything in this chapter is a hypothesis with a
named test. None is yet a finding.

\section{The two ends are different physics}

The author puts the distinction sharply: a sick person's immune cells working hard to make repairs does not mean their informational fidelity is compromised;
and an elevated reading is not the same event as a severely suppressed one.

\begin{description}
\item[Suppressed ($A<0.95$).] The cell holds less entropy on its identity sites than the healthy cell of its kind. In the author's reading the writing process has
      slowed or stopped and identity maintenance is failing: senescence, exhaustion, a lineage running down.
\item[Elevated ($A>1.05$).] More entropy than the healthy cell of its kind. In the author's image, the printer is in overdrive: more is being
      written than the surface can hold as a pattern, and the surface saturates. This is the direction of the horizon analogy of Chapter~\ref{ch:floorbreach}.
\item[Activity is neither.] A neutrophil clearing an infection does more work while keeping its identity. The physics does not predict that a busy immune cell
      moves on the gauge, and nothing on the report may say that it does.
\end{description}
All three are \conjecture.

\section{Four mechanisms, four tests}

\begin{center}\small
\begin{longtable}{@{}p{0.2\textwidth}p{0.16\textwidth}p{0.28\textwidth}p{0.26\textwidth}@{}}
\toprule mechanism & end & which cells & test \\\midrule\endhead
the leukocyte is the tumour: leukaemias, lymphoma, myeloma; clonal haematopoiesis as the pre-malignant form & elevated & the affected lineage; the others near 1.00. The diseased cell is the majority cell, so the fraction limit cannot fake it & per-cell $A$ on public AML and CLL arrays, affected lineage against the others, with held-out healthy progenitors \\
ageing & elevated & slowly, per person, over years & serial mode on SATSA draws over 10--19 years \\
T-cell exhaustion, immunosenescence & suppressed & CD4 and CD8 specifically, with neutrophils and monocytes near 1.00 in the same specimen & per-cell direction test on Alzheimer's and progressive supranuclear palsy arrays already collected \\
emergency haematopoiesis: immature cells pushed out of the marrow & composition first & progenitor fraction up; immature cells may read suppressed & composition and progenitor readings on the same arrays \\\bottomrule
\end{longtable}
\end{center}

\section{Ageing is on the elevated side}

The first chain's record and the first readings on the new atlas agree in direction: old immune cells read higher than young ones. The oldest Uppsala arrays read
immune $A$ about 1.037 on the first atlas, and neutrophil $A$ rose with age across 24 arrays on the second (Chapter~\ref{ch:firstreadings}). \observed{} The author's
reading is that an old immune system has had more written over a lifetime onto the same surface, with less headroom, which is the same direction as clonal
haematopoiesis. Whether old blood's elevation is a slope shared by every cell or a clone riding on it is exactly what per-cell trajectories separate: a slope moves
every lineage together; a clone moves one lineage away from the others.

\section{What the leukocyte gauge does not do}

It does not detect a solid tumour through the blood cells' $A$. Immune cells do not take on a tumour's state by fighting it. A solid cancer, if it shows in whole blood
at all, shows through composition or through shed cells, and on whole-blood arrays shed cells below about 2~\% of the DNA are below what the array can see.
\measured{} Plasma cell-free DNA, where shed cells are a much larger share, is the specimen for that question (Chapter~\ref{part4:ch:reach}).

\section{Where the ``bidirectional'' picture came from}

Earlier engines found the immune metric moving in opposite directions by condition. Those were pooled scores over many loci and many cells, computed as a case group
minus a control group. Two things produce opposite signs there without any cell going both ways. Pooling: at the locus level some CpGs drift up and others down and
the entropy of the mean over all of them nets to nothing; at the cell level one lineage moves one way and another the other, netted by a composition that itself moves
with disease and age. And the comparison: a control group sitting above or below 1.00 flips the sign of a departure. Neither exists in the present instrument. Each cell
is read on its own identity loci, chosen so the statistic is unimodal, against 1.00. If T cells read 1.06 and B cells 0.97 in one specimen, the report prints two rows
and nothing cancels.

\section{A count, not a score}

For conditions that involve the immune system as a whole, the instrument reports how many immune cells read outside Normal, and which. That is a count of cell
readings, not a reading of the class, and it is the only form in which a class appears on a report as a group.
